\color{unidarkgreen}{\subsection{Логика отключения}}
\color{black}

\begin{enumerate}[label=\arabic{section}.\arabic{subsection}.\arabic*,  align=left, labelsep=4pt, leftmargin=0pt, itemindent=57pt]

\item
Устройство формирует выходное воздействие на отключение В1 ВН, В2 ВН и ТМПд по схеме
«ИЛИ» в следующих случаях:
\begin{itemize}
\item Срабатывание ДТЗ;
\item Срабатывание МТЗ-2 КО;
\item Срабатывание ТЗНП-2 КО;
\item Срабатывание ТЗОП-2 КО;
\item Срабатывание МТЗ-2 НН;
\item Срабатывание МТЗ ОУ;
\item Срабатывание КИ НН;
\item Отключение от ЗДЗ НН;
\item Отключение от УРОВ ТМПо;
\item Отключение от УРОВ ТМПр;
\item Отключение от внешней РЗ КО, ОУ.
\end{itemize}

\item
Устройство формирует выходное воздействие на отключение ТМПо, ТМПр по схеме «ИЛИ» в
следующих случаях:
\begin{itemize}
\item Срабатывание МТЗ-1 КО;
\item Срабатывание ТЗНП-1 КО;
\item Срабатывание ТЗОП-1 КО;
\item Срабатывание МТЗ-1 НН.
\end{itemize}

\item
Сброс фиксации производится с помощь \textit{«Сброс»}.

\begin{figure}[H]
\centering
\includegraphics[width=0.55\textwidth,height=0.55\textheight,keepaspectratio]{\fbpath/Логика отключения/img/LO.pdf}
\captionof{figure}{Функциональный блок логики отключения}\label{pict_lo:LO}
\end{figure}

\end{enumerate}